\documentclass{article}

\usepackage{amsthm}
\usepackage{keytheorems}
\usepackage{hyperref} 

%Interface théorème
\newkeytheorem{lem}[
	name = Lemme
]

\newkeytheorem{prop}[
	name = Proposition
]

\newkeytheorem{thm}[
	name = Théorème
]

\newkeytheorem{coro}[
	name = Corollaire
]

\renewcommand*{\proofname}{Démonstration}

\theoremstyle{definition}
\newtheorem{defn}{Définition}

\theoremstyle{definition}
\newtheorem*{nota}{Notation}

\theoremstyle{definition}
\newtheorem*{limi}{Liminaire}

\theoremstyle{remark}
\newtheorem{rmq}{Remarque}

\theoremstyle{plain}
\newtheorem*{fait}{Fait}

\newcommand{\V}{\text{V}}
\newcommand{\Vt}{\tilde{V}}

\begin{document}

\begin{thm}
Soient $v_1, \dots, v_n$ $n$ vecteurs de $\R^n$ ; on note $\V(v_1, \dots, v_n)$ le volume du parallélépipède engengré par les vecteurs 
$v_1, \dots, v_n$. On a alors :
\[ \V(v_1, \dots, v_n) = \abs{\det (v_1, \dots, v_n)} \]
\end{thm}

\begin{proof}
Posons :

\[ \Vt(v_1, \dots, v_n) = \left \{ 
	\begin{array}{c c c}  
		\V(v_1, \dots, v_n) & \quad \text{si} \quad & \det (v_1, \dots, v_n) \geq 0 \\
		- \V(v_1, \dots, v_n) & \quad \text{si} \quad & \det (v_1, \dots, v_n) < 0
	\end{array} \right.
\]

On va montrer que $\Vt(v_1, \dots, v_n) = \det (v_1, \dots, v_n)$ ; pour cela, il suffit de montrer que la forme linéaire $\Vt$ est 
$n$-linéaire, alternée et telle que $\Vt(e_1, \dots, e_n) = 1$ avec $(e_1, \dots, e_n)$ la base canonique de $K^n$. \newline

Si la famille $(v_1, \dots, v_n)$ est liée, on a immédiatement $\Vt(v_1, \dots, v_n) = 0$, ce qui démontrer que $\Vt$ est alternée. \newline

Il reste à démontrer la $n$-linéarité de $\Vt$. Démontrons d'abord que :
\[ \forall \lmabda \in \R, \forall i \in \{1, \dots, n \} ; \V(v_1, \dots, \lambda v_i, \dots, v_n) = \abs{\lambda) \V(v_1, \dots, v_n) \]
Le résultat est évident pour $\lambda = 0$, ou si la famille $(v_1, \dots, v_n)$ est liée. Pour $\lambda=-1

\end{proof}

\end{document}